% Options for packages loaded elsewhere
\PassOptionsToPackage{unicode}{hyperref}
\PassOptionsToPackage{hyphens}{url}
\documentclass[
  ignorenonframetext,
]{beamer}
\newif\ifbibliography
\usepackage{pgfpages}
\setbeamertemplate{caption}[numbered]
\setbeamertemplate{caption label separator}{: }
\setbeamercolor{caption name}{fg=normal text.fg}
\beamertemplatenavigationsymbolsempty
% remove section numbering
\setbeamertemplate{part page}{
  \centering
  \begin{beamercolorbox}[sep=16pt,center]{part title}
    \usebeamerfont{part title}\insertpart\par
  \end{beamercolorbox}
}
\setbeamertemplate{section page}{
  \centering
  \begin{beamercolorbox}[sep=12pt,center]{section title}
    \usebeamerfont{section title}\insertsection\par
  \end{beamercolorbox}
}
\setbeamertemplate{subsection page}{
  \centering
  \begin{beamercolorbox}[sep=8pt,center]{subsection title}
    \usebeamerfont{subsection title}\insertsubsection\par
  \end{beamercolorbox}
}
% Prevent slide breaks in the middle of a paragraph
\widowpenalties 1 10000
\raggedbottom
\AtBeginPart{
  \frame{\partpage}
}
\AtBeginSection{
  \ifbibliography
  \else
    \frame{\sectionpage}
  \fi
}
\AtBeginSubsection{
  \frame{\subsectionpage}
}
\usepackage{iftex}
\ifPDFTeX
  \usepackage[T1]{fontenc}
  \usepackage[utf8]{inputenc}
  \usepackage{textcomp} % provide euro and other symbols
\else % if luatex or xetex
  \usepackage{unicode-math} % this also loads fontspec
  \defaultfontfeatures{Scale=MatchLowercase}
  \defaultfontfeatures[\rmfamily]{Ligatures=TeX,Scale=1}
\fi
\usepackage{lmodern}
\ifPDFTeX\else
  % xetex/luatex font selection
\fi
% Use upquote if available, for straight quotes in verbatim environments
\IfFileExists{upquote.sty}{\usepackage{upquote}}{}
\IfFileExists{microtype.sty}{% use microtype if available
  \usepackage[]{microtype}
  \UseMicrotypeSet[protrusion]{basicmath} % disable protrusion for tt fonts
}{}
\makeatletter
\@ifundefined{KOMAClassName}{% if non-KOMA class
  \IfFileExists{parskip.sty}{%
    \usepackage{parskip}
  }{% else
    \setlength{\parindent}{0pt}
    \setlength{\parskip}{6pt plus 2pt minus 1pt}}
}{% if KOMA class
  \KOMAoptions{parskip=half}}
\makeatother
\setlength{\emergencystretch}{3em} % prevent overfull lines
\providecommand{\tightlist}{%
  \setlength{\itemsep}{0pt}\setlength{\parskip}{0pt}}
% Auto-generated by infrastructure/rendering/_pdf_latex_helpers.py::extract_math_font_preamble.
% Minimal header passed to Pandoc via -H so Beamer slides pick up the same math font
% as the combined-PDF path. Adjust the manuscript's preamble.md to change the font globally.
\usepackage{unicode-math}
\setmathfont{latinmodern-math.otf}

% Slide layout defaults for warning-clean scientific decks.
\usepackage{etoolbox}
\IfFileExists{xurl.sty}{\usepackage{xurl}}{}
\IfFileExists{seqsplit.sty}{\usepackage{seqsplit}}{\newcommand{\seqsplit}[1]{#1}}
\protected\def\breakseq#1{\seqsplit{#1}}
\protected\def\breaktt#1{\begingroup\ttfamily\seqsplit{#1}\endgroup}
\setlength{\emergencystretch}{6em}
\tolerance=5000
\hbadness=10000
\hfuzz=1pt
\setlength{\tabcolsep}{2pt}
\AtBeginEnvironment{longtable}{\tiny\renewcommand{\arraystretch}{0.86}\setlength{\tabcolsep}{1pt}}
\AtBeginEnvironment{tabular}{\tiny\renewcommand{\arraystretch}{0.86}\setlength{\tabcolsep}{1pt}}

% Natbib and cross-reference fallbacks — slides don't load natbib
% or cleveref, but combined-PDF manuscript prose may emit these
% commands. The fallback renders citations as a bracketed key list
% and cross-references as detokenized labels so slides don't fail on
% undefined control sequences or raw underscores. \providecommand is
% a no-op if packages load later.
\providecommand{\citep}[1]{[#1]}
\providecommand{\citet}[1]{#1}
\providecommand{\citealp}[1]{#1}
\providecommand{\citeauthor}[1]{#1}
\providecommand{\citeyear}[1]{#1}
\providecommand{\cref}[1]{\texttt{\detokenize{#1}}}
\providecommand{\Cref}[1]{\texttt{\detokenize{#1}}}
\usepackage{bookmark}
\IfFileExists{xurl.sty}{\usepackage{xurl}}{} % add URL line breaks if available
\urlstyle{same}
% fallback for those not using the hyperref driver hyperxmp:
\makeatletter
\@ifundefined{xmpquote}{\newcommand{\xmpquote}[1]{#1}}{}
\makeatother
\hypersetup{
  hidelinks,
  pdfcreator={LaTeX via pandoc}}

\author{\texorpdfstring{}{}}
\date{}

\begin{document}

\section{Introduction}\label{sec:introduction}

\begin{frame}[fragile,allowframebreaks]{Introduction}
This \breaktt{template\_eda\_notebook} serves as the
exploratory-data-analysis exemplar for the
\href{https://github.com/docxology/template}{Research Project Template}
ecosystem, demonstrating how a computational notebook can stay
reproducible by delegating every computation to a fully-tested library.
The prose, the labelled figures, and the summary table are all produced
through an auditable custody chain: a deterministic dataset, tested
functions in \texttt{src/eda/}, a thin analysis script, and multi-format
rendering.
\end{frame}

\begin{frame}[fragile,allowframebreaks]{Why exploratory data analysis}
\protect\phantomsection\label{why-exploratory-data-analysis}
EDA --- the work of getting to know a dataset before committing to a
model --- is where most research projects actually begin: load the data,
see how much is missing, look at distributions, and check which
variables move together. It is fast and interactive by nature, which is
exactly why it tends to live in notebooks. The hazard is that the
interactive convenience of a notebook cell is also a trap: a one-off
\breaktt{df.groupby(...).mean()} becomes load-bearing, never gets a test,
and silently disagrees with the figure two cells down after the data
changes.
\end{frame}

\begin{frame}[fragile,allowframebreaks]{The notebook -\textgreater{}
tested src extraction workflow}
\protect\phantomsection\label{the-notebook---tested-src-extraction-workflow}
This exemplar teaches one discipline: \textbf{explore in a cell, then
extract to a tested library the moment a computation matters.}
Concretely:

\begin{enumerate}
\tightlist
\item
  The walkthrough notebook (\breaktt{notebooks/eda\_walkthrough.ipynb})
  imports from \texttt{src} and calls tested functions; it contains no
  business logic of its own.
\item
  Each analytical step --- loading, cleaning, summarizing, correlating,
  preparing figure data --- is a typed, documented function in
  \texttt{src/eda/} with a test that asserts an exact numeric property.
\item
  A thin script (\breaktt{scripts/eda\_analysis.py}) runs the same
  pipeline headless and writes the figures and a summary CSV to
  \texttt{output/}.
\end{enumerate}
\end{frame}

\begin{frame}[fragile,allowframebreaks]{Template architecture context}
\protect\phantomsection\label{template-architecture-context}
The project sits on the repository's three pillars:

\begin{enumerate}
\tightlist
\item
  \textbf{\texttt{src/eda/} library}: pure pandas/numpy data transforms
  --- no plotting, no file I/O, no \texttt{infrastructure} imports. This
  purity is what makes the library forkable and trivially testable.
\item
  \textbf{\texttt{tests/} framework}: a zero-mock suite that exercises
  the library against the shipped CSV and tiny real frames, plus a
  structural check that the notebook's imports bind to the library's
  public surface.
\item
  \textbf{\texttt{docs/} knowledge base}: architectural guidelines, the
  testing philosophy, and the operational rules that govern agents
  editing this tree.
\end{enumerate}
\end{frame}

\begin{frame}[allowframebreaks]{The dataset}
\protect\phantomsection\label{the-dataset}
We analyze a small synthetic cohort of subject measurements --- height
(cm), weight (kg), and resting heart rate (bpm) across three groups ---
generated with a fixed seed so every statistic in \breakseq{[@sec:results]} is
reproducible. The data is shaped so that weight depends positively on
height (a strong, easy-to-see correlation) while resting heart rate is
only weakly related, and a few cells are left blank to exercise the
missing-data path honestly.
\end{frame}

\begin{frame}[fragile,allowframebreaks]{Reader's guide to the
manuscript}
\protect\phantomsection\label{readers-guide-to-the-manuscript}
\begin{itemize}
\tightlist
\item
  \textbf{\breakseq{[@sec:methodology]}} ties each EDA step to its function in
  \texttt{src/eda/}.
\item
  \textbf{\breakseq{[@sec:results]}} is figure-centric: each panel names the
  figure-data preparer that produced it.
\item
  \textbf{\breakseq{[@sec:experimental\_setup]}} lists the dataset schema and
  software environment.
\item
  \textbf{\breakseq{[@sec:reproducibility]}} records the artifact inventory and
  the exact commands to regenerate everything.
\item
  \textbf{\breakseq{[@sec:scope]}} states scope and related literature so the
  exemplar is not mistaken for a general-purpose EDA toolkit.
\end{itemize}
\end{frame}

\end{document}
